\section{Hasil Pengujian}

Pengujian sistem dilakukan untuk mengetahui apakah sistem yang telah dibuat dapat berjalan sesuai dengan yang diharapkan. Pengujian ini dilakukan dengan cara melakukan uji coba terhadap setiap fitur yang ada pada sistem, baik dari sisi \textit{User} maupun dari sisi admin.

\subsection{Pengujian Blackbox}

Pengujian Blackbox adalah pengujian yang dilakukan untuk mengetahui apakah sistem yang telah dibuat dapat berjalan sesuai dengan yang diharapkan. Pengujian ini dilakukan dengan cara melakukan uji coba terhadap setiap fitur yang ada pada sistem, baik dari sisi \textit{User} pada tabel \ref{tabel:hasil_bb_user}, \textit{Admin} pada tabel \ref{tabel:uat_admin_cabang_hasil}, dan \textit{SuperAdmin} pada tabel \ref{tabel:hasil_bb_superadmin}. Pengujian ini dilakukan tanpa melihat kode program yang ada pada sistem dan hanya berfokus pada input dan output dari interaksi antara pengguna dengan sistem.

\begin{longtable}{|p{3.8cm}|p{4.5cm}|p{4.5cm}|p{2cm}|}
    \caption{Hasil Pengujian Black Box (User)} \label{tabel:hasil_bb_user}                                                                                                                                                                                                                                                                                                                                                                                                                                                                                                   \\
    \hline
    \textbf{Skenario Pengujian}                                                & \textbf{Langkah Pengujian}                                                                                                                                                                                                                                                                                        & \textbf{Hasil yang diharapkan}                                                                                                              & \textbf{Kesesuaian Hasil} \\
    \hline
    \endfirsthead

    % \multicolumn{4}{c}%
    % {{\bfseries Tabel \thetable\ lanjut dari halaman sebelumnya}}                                                                                                                                                                                                                                                                                                                                                                                                                                                                                                            \\
    % \hline
    % \textbf{Skenario Pengujian}                                                & \textbf{Langkah Pengujian}                                                                                                                                                                                                                                                                                        & \textbf{Hasil yang diharapkan}                                                                                                              & \textbf{Kesesuaian Hasil} \\
    % \hline
    % \endhead

    % \hline \multicolumn{4}{r}{\textit{Bersambung ke halaman berikutnya}}                                                                                                                                                                                                                                                                                                                                                                                                                                                                                                     \\
    % \endfoot

    \hline
    \endlastfoot

    % --- Konten tabel Anda dimulai di sini ---
    \multicolumn{4}{|l|}{\textbf{Fitur Registrasi Akun}}
    \\
    \hline
    Registrasi berhasil dengan email baru                                      & 1. Buka halaman registrasi. \newline 2. Isi nama lengkap. \newline 3. Masukkan email yang belum terdaftar (format valid). \newline 4. Masukkan kata sandi kuat. \newline 5. Konfirmasi kata sandi dengan benar. \newline 6. (Jika ada) Centang persetujuan Syarat \& Ketentuan. \newline 7. Klik tombol "Daftar". & - Pesan sukses registrasi ditampilkan. \newline - Pengguna diarahkan ke halaman login atau dashboard. \newline - Email konfirmasi terkirim. & Sesuai                    \\
    \hline
    Registrasi gagal karena email sudah terdaftar                              & 1. Buka halaman registrasi. \newline 2. Isi nama lengkap. \newline 3. Masukkan email yang sudah terdaftar. \newline 4. Masukkan kata sandi. \newline 5. Konfirmasi kata sandi. \newline 6. Klik tombol "Daftar".                                                                                                  & - Pesan error "Email sudah terdaftar" ditampilkan. \newline - Pengguna tetap di halaman registrasi.                                         & Sesuai                    \\
    \hline
    Registrasi gagal karena format email salah                                 & 1. Buka halaman registrasi. \newline 2. Masukkan email dengan format salah (misal: user@.com, user.com). \newline 3. Isi field lainnya. \newline 4. Klik "Daftar".                                                                                                                                                & - Pesan error "Format email tidak valid" ditampilkan.                                                                                       & Sesuai                    \\
    \hline
    Registrasi gagal karena konfirmasi kata sandi tidak cocok                  & 1. Buka halaman registrasi. \newline 2. Isi nama lengkap, email baru, kata sandi. \newline 3. Masukkan konfirmasi kata sandi yang berbeda. \newline 4. Klik tombol "Daftar".                                                                                                                                      & - Pesan error "Konfirmasi kata sandi tidak cocok" ditampilkan.                                                                              & Sesuai                    \\
    \hline
    Registrasi menggunakan akun Google (sukses, pengguna baru)                 & 1. Buka halaman registrasi/login. \newline 2. Klik tombol "Daftar/Masuk dengan Google". \newline 3. Pilih akun Google yang valid dan belum terdaftar di sistem. \newline 4. Izinkan akses jika diminta.                                                                                                           & - Akun berhasil terbuat/terhubung. \newline - Pengguna otomatis berhasil masuk.                                                             & Sesuai                    \\
    \hline

    \multicolumn{4}{|l|}{\textbf{Fitur Login Pengguna}}
    \\
    \hline
    Login berhasil dengan email dan kata sandi valid                           & 1. Buka halaman login. \newline 2. Masukkan email terdaftar. \newline 3. Masukkan kata sandi yang benar. \newline 4. Klik tombol "Masuk".                                                                                                                                                                         & - Login berhasil. \newline - Pengguna berhasil masuk.                                                                                       & Sesuai                    \\
    \hline
    Login gagal (email tidak terdaftar)                                        & 1. Buka halaman login. \newline 2. Masukkan email yang tidak terdaftar. \newline 3. Masukkan kata sandi. \newline 4. Klik tombol "Masuk".                                                                                                                                                                         & - Pesan error "Email atau kata sandi salah" ditampilkan.                                                                                    & Sesuai                    \\
    \hline
    Login gagal (kata sandi salah)                                             & 1. Buka halaman login. \newline 2. Masukkan email terdaftar. \newline 3. Masukkan kata sandi yang salah. \newline 4. Klik tombol "Masuk".                                                                                                                                                                         & - Pesan error "Email atau kata sandi salah" ditampilkan.                                                                                    & Sesuai                    \\
    \hline
    Login gagal (akun belum diverifikasi email, jika ada)                      & 1. Registrasi dengan email lalu jangan verifikasi email. \newline 2. Coba login dengan email dan kata sandi tersebut.                                                                                                                                                                                             & - Pesan error "Akun belum diverifikasi. Silakan cek email Anda." atau sejenisnya.                                                           & Sesuai                    \\
    \hline
    Login menggunakan akun Google (sukses, akun sudah ada)                     & 1. Buka halaman login. \newline 2. Klik "Masuk dengan Google". \newline 3. Pilih akun Google yang sudah terdaftar/terhubung sebelumnya.                                                                                                                                                                           & - Login berhasil. \newline - Pengguna diarahkan ke dashboard.                                                                               & Sesuai                    \\
    \hline

    \multicolumn{4}{|l|}{\textbf{Fitur Lupa Kata Sandi}}
    \\
    \hline
    Proses lupa kata sandi berhasil                                            & 1. Di halaman login, klik "Lupa Kata Sandi?". \newline 2. Masukkan email terdaftar. \newline 3. Klik "Kirim Instruksi". \newline 4. Buka email, klik link reset. \newline 5. Masukkan kata sandi baru dan konfirmasinya. \newline 6. Klik "Reset Kata Sandi".                                                     & - Pesan sukses "Kata sandi berhasil direset". \newline - Dapat login dengan kata sandi baru.                                                & Sesuai                    \\
    \hline
    Lupa kata sandi dengan email tidak terdaftar                               & 1. Di halaman login, klik "Lupa Kata Sandi?". \newline 2. Masukkan email yang tidak terdaftar. \newline 3. Klik "Kirim Instruksi".                                                                                                                                                                                & - Pesan error "Email tidak ditemukan" atau instruksi tidak terkirim.                                                                        & Sesuai                    \\
    \hline

    \multicolumn{4}{|l|}{\textbf{Fitur Kelola Profil Pengguna}}
    \\
    \hline
    Mengubah semua data profil yang diizinkan (sukses)                         & 1. Login. \newline 2. Ke halaman "Profil Saya" lalu "Edit Profil". \newline 3. Ubah nama, nomor telepon, alamat, gender, tanggal lahir. \newline 4. Klik "Simpan".                                                                                                                                                & - Pesan sukses "Profil berhasil diperbarui". \newline - Semua perubahan tercermin dengan benar.                                             & Sesuai                    \\
    \hline
    Mengubah data profil dengan format salah (misal: telepon)                  & 1. Login. \newline 2. Ke "Edit Profil". \newline 3. Masukkan nomor telepon dengan format salah (misal: huruf). \newline 4. Klik "Simpan".                                                                                                                                                                         & - Pesan error "Format nomor telepon tidak valid". \newline - Data tidak tersimpan.                                                          & Sesuai                    \\
    \hline

    \multicolumn{4}{|l|}{\textbf{Fitur Verifikasi Identitas (KYC)}}
    \\
    \hline
    Mengajukan verifikasi KYC (KTP) dengan data valid                          & 1. Login. \newline 2. Ke halaman KYC. \newline 3. Pilih jenis ID "KTP". \newline 4. Masukkan nomor KTP valid. \newline 5. Unggah foto KTP jelas. \newline 6. Unggah foto selfie jelas. \newline 7. Klik "Ajukan".                                                                                                 & - Pengajuan KYC berhasil dikirim. \newline - Status KYC menjadi "Menunggu Verifikasi".                                                      & Sesuai                    \\
    \hline
    Mengajukan KYC (Passport) dengan data valid                                & 1. Login. \newline 2. Ke halaman KYC. \newline 3. Pilih jenis ID "Passport". \newline 4. Masukkan nomor Passport valid. \newline 5. Unggah foto Passport jelas. \newline 6. Unggah foto selfie jelas. \newline 7. Klik "Ajukan".                                                                                  & - Pengajuan KYC berhasil dikirim. \newline - Status KYC menjadi "Menunggu Verifikasi".                                                      & Sesuai                    \\
    \hline
    Mengajukan KYC dengan file gambar terlalu besar/format salah               & 1. Login. \newline 2. Ke halaman KYC. \newline 3. Isi data. \newline 4. Unggah file gambar dengan ukuran > batas maksimal atau format bukan JPG/PNG. \newline 5. Klik "Ajukan".                                                                                                                                   & - Pesan error "Ukuran file terlalu besar" atau "Format file tidak didukung".                                                                & Sesuai                    \\
    \hline
    Melihat status KYC (Pending, Approved, Rejected)                           & 1. Login. \newline 2. Navigasi ke halaman KYC atau Profil.                                                                                                                                                                                                                                                        & - Status KYC pengguna ditampilkan dengan benar sesuai kondisi aktual.                                                                       & Sesuai                    \\
    \hline

    \multicolumn{4}{|l|}{\textbf{Fitur Pencarian dan Reservasi}}
    \\
    \hline
    Mencari pod (tanggal check-out sebelum check-in)                           & 1. Login. \newline 2. Pilih Cabang. \newline 3. Pilih tanggal check-in. \newline 4. Pilih tanggal check-out lebih awal dari check-in. \newline 5. Klik "Cari Pod".                                                                                                                                                & - Pesan error "Tanggal check-out tidak valid".                                                                                              & Sesuai                    \\
    \hline
    Melakukan pemesanan pod dengan voucher kadaluarsa                          & 1. Lakukan pencarian pod. \newline 2. Pilih tipe pod. \newline 3. Di halaman pembayaran, masukkan kode voucher yang sudah kadaluarsa. \newline 4. Klik "Terapkan Voucher".                                                                                                                                        & - Pesan error "Voucher tidak valid atau sudah kadaluarsa".                                                                                  & Sesuai                    \\
    \hline
    Pemesanan pod dengan jumlah melebihi kapasitas tipe pod                    & 1. Lakukan pencarian pod. \newline 2. Pilih tipe pod. \newline 3. Coba pesan jumlah tamu yang melebihi kapasitas maksimal dewasa/anak tipe pod tersebut.                                                                                                                                                          & - Pesan error "Jumlah tamu melebihi kapasitas" atau input dibatasi.                                                                         & Sesuai                    \\
    \hline
    Pemesanan pod, pembayaran berhasil (QRIS)                                  & 1. Lakukan pencarian pod. \newline 2. Pilih tipe pod. \newline 3. Lanjutkan ke pembayaran. \newline 4. Pilih metode QRIS, pindai QR. \newline 5. Konfirmasi pembayaran di aplikasi e-wallet.                                                                                                                      & - Pembayaran berhasil. \newline - Tampilan menunggu dibayar berupah menjadi terbayar \newline - Transaksi tercatat.                         & Sesuai                    \\
    \hline
    Pemesanan pod, menggunakan kode referral teman (valid, aturan disesuaikan) & 1. Lakukan pencarian pod. \newline 2. Pilih tipe pod. \newline 3. Di halaman pembayaran, masukkan kode referral valid milik teman. \newline 4. Klik "Terapkan". \newline 5. Selesaikan pembayaran.                                                                                                                & - Diskon referral berhasil diterapkan. \newline - Harga total berkurang. \newline - Pemesanan berhasil.                                     & Sesuai                    \\
    \hline

    \multicolumn{4}{|l|}{\textbf{Fitur Sinkronisasi OTA}}
    \\
    \hline
    Sinkronisasi reservasi OTA dengan kode valid                               & 1. Login. \newline 2. Ke halaman "Sinkronisasi OTA" atau "Klaim Reservasi". \newline 3. Pilih OTA. \newline 4. Masukkan kode booking OTA / nomor order OTA yang valid (sudah diinput Admin Cabang). \newline 5. Klik "Sinkronisasi".                                                                              & - Reservasi OTA berhasil disinkronkan/diklaim. \newline - Muncul di daftar booking pengguna.                                                & Sesuai                    \\
    \hline
    Sinkronisasi reservasi OTA dengan kode tidak ditemukan/salah               & 1. Login. \newline 2. Ke halaman "Sinkronisasi OTA". \newline 3. Masukkan kode booking OTA yang salah. \newline 4. Klik "Sinkronisasi".                                                                                                                                                                           & - Pesan error "Reservasi tidak ditemukan" atau "Kode booking salah".                                                                        & Sesuai                    \\
    \hline

    \multicolumn{4}{|l|}{\textbf{Fitur Operasional Menginap}}
    \\
    \hline
    Check-in berhasil                                                          & 1. Login. \newline 2. Buka halaman stays lalu klik tab current \newline 3. Coba klik tombol "Check-in" \newline 4. Pilih pod yang diinginkan \newline 5. klik checkin                                                                                                                                             & - berhasil check-in dan mendapatkan kunci pod yang diinginkan berupa QRCode                                                                 & Sesuai                    \\
    \hline
    Check-in gagal (di luar jam operasional atau sebelum waktunya check-in)    & 1. Login. \newline 2. Buka halaman stays. \newline 3. Coba klik tombol "Check-in" sebelum jam check-in yang diizinkan atau setelah batas waktu check-in.                                                                                                                                                          & - Tombol "Check-in" tidak ada untuk booking yang stausnya masih upcoming                                                                    & Sesuai                    \\
    \hline
    Check-out                                                                  & 1. Login. \newline 2. Buka detail booking aktif yang sudah check-in. \newline 3. Klik tombol "Check-out" \newline 4. muncul modal rating (user wajib mengisi rating untuk isi review opsional) \newline 5. klik kirim                                                                                             & - Check-out berhasil. \newline - Status booking "Completed". \newline rating dan review juga berhasil dikirim                               & Sesuai                    \\
    \hline

    \multicolumn{4}{|l|}{\textbf{Fitur Referral}}
    \\
    \hline
    Melihat riwayat pendapatan referral                                        & 1. Login. \newline 2. Ke halaman Referral. \newline 3. Pilih tab "Riwayat Pendapatan".                                                                                                                                                                                                                            & - Menampilkan daftar transaksi yang menghasilkan pendapatan referral dengan statusnya.                                                      & Sesuai                    \\
    \hline
    Mengajukan pencairan referral (KYC belum disetujui)                        & 1. Login, KYC masih "Pending" atau "Rejected". \newline 2. Ke halaman Referral. \newline 3. Coba klik "Cairkan Dana".                                                                                                                                                                                             & - Tombol "Cairkan Dana" tidak aktif atau pesan error "Verifikasi KYC diperlukan untuk pencairan".                                           & Sesuai                    \\
    \hline

    \multicolumn{4}{|l|}{\textbf{Fitur Ulasan}}
    \\
    \hline
    Memberikan ulasan tanpa rating bintang                                     & 1. Login, setelah checkout. \newline 2. Ke Riwayat Booking, pilih tab oldest booking. \newline 3. Klik "Beri Ulasan". \newline 4. Hanya tulis komentar, tidak memilih bintang. \newline 5. Coba Kirim.                                                                                                            & - tombol kirim review tidak dapat di klik                                                                                                   & Sesuai                    \\
    \hline
    Mencoba memberi ulasan untuk booking yang belum selesai/dibatalkan         & 1. Login. \newline 2. Cari booking yang statusnya "Pending Payment" atau "Cancelled". \newline 3. Coba cari opsi "Beri Ulasan".                                                                                                                                                                                   & - Opsi "Beri Ulasan" tidak tersedia untuk booking tersebut.                                                                                 & Sesuai                    \\
    \hline
\end{longtable}

\begin{longtable}{|p{3.8cm}|p{4.5cm}|p{4.5cm}|p{2cm}|} % Menggunakan format 4 kolom
    \caption{Hasil Pengujian Black Box (Admin)} \label{tabel:hasil_bb_admin}                                                                                                                                                                                                                                                                                                                                                                                                                    \\
    \hline
    \textbf{Skenario Pengujian}                                     & \textbf{Langkah Pengujian}                                                                                                                                                                                                                          & \textbf{Hasil yang diharapkan}                                                                                                          & \textbf{Kesesuaian Hasil} \\
    \hline
    \endfirsthead

    % \multicolumn{4}{c}%
    % {{\bfseries Tabel \thetable\ lanjut dari halaman sebelumnya}}                                                                                                                                                                                                                                                                                                                                                                                                                               \\
    % \hline
    % \textbf{Skenario Pengujian}                                     & \textbf{Langkah Pengujian}                                                                                                                                                                                                                          & \textbf{Hasil yang diharapkan}                                                                                                          & \textbf{Kesesuaian Hasil} \\
    % \hline
    % \endhead

    % \hline \multicolumn{4}{r}{\textit{Bersambung ke halaman berikutnya}}                                                                                                                                                                                                                                                                                                                                                                                                                        \\
    % \endfoot

    \hline
    \endlastfoot

    % --- Konten tabel Admin Cabang ---
    \multicolumn{4}{|l|}{\textbf{Fitur Login dan Profil Admin Cabang}}
    \\
    \hline
    Login Admin Cabang berhasil                                     & 1. Buka halaman login admin. \newline 2. Masukkan email Admin Cabang valid. \newline 3. Masukkan kata sandi benar. \newline 4. Klik "Masuk".                                                                                                        & - Login berhasil. \newline - Diarahkan ke halaman pilih cabang yang ingin dikelolan Cabang.                                             & Sesuai                    \\
    \hline
    Login Admin Cabang gagal (kata sandi salah)                     & 1. Buka halaman login admin. \newline 2. Masukkan email Admin Cabang valid. \newline 3. Masukkan kata sandi salah. \newline 4. Klik "Masuk".                                                                                                        & - Pesan error "Email atau kata sandi salah". \newline - Tetap di halaman login.                                                         & Sesuai                    \\
    \hline
    Mengubah profil Admin Cabang (misal: nama)                      & 1. Login sebagai Admin Cabang. \newline 2. Navigasi ke "Profil Saya". \newline 3. Klik "Edit Profil". \newline 4. Ubah nama. \newline 5. Klik "Simpan".                                                                                             & - Nama Admin Cabang berhasil diperbarui.                                                                                                & Sesuai                    \\
    \hline
    \multicolumn{4}{|l|}{\textbf{Fitur Pengelolaan Operasional Cabang}}
    \\
    \hline
    Mengedit semua informasi detail cabang yang diizinkan           & 1. Login. \newline 2. Ke "Pengaturan data Cabang". \newline 3. Ubah nama, alamat, telepon, email, deskripsi, jam check-in/out, gambar utama. \newline 4. Klik "Simpan".                                                                             & - Semua perubahan berhasil disimpan dan ditampilkan dengan benar.                                                                       & Sesuai                    \\
    \hline
    Menambah tipe pod baru dengan semua field valid                 & 1. Login. \newline 2. Ke "Menu jenis pod" \newline 3. Klik "Tambah Baru". \newline 4. Isi nama, deskripsi, harga, kapasitas (dewasa, anak), gambar, fasilitas terkait, status breakfast. \newline 5. Klik "Simpan".                                 & - Tipe pod baru berhasil ditambahkan dengan semua detailnya.                                                                            & Sesuai                    \\
    \hline
    Menghapus tipe pod yang masih memiliki pod aktif                & 1. Login. \newline 2. Ke "Manajemen Pod" - "Tipe Pod". \newline 3. Pilih tipe pod yang memiliki pod terdaftar di bawahnya. \newline 4. Klik "Hapus". \newline 5. Konfirmasi.                                                                        & - Pesan error "Tidak dapat menghapus tipe pod karena masih ada pod yang terhubung" atau sejenisnya. \newline - Tipe pod tidak terhapus. & Sesuai                    \\
    \hline
    Menambah pod baru dengan device code IoT unik                   & 1. Login. \newline 2. Ke "Manajemen Pod" - "pilih jenis pod yang pod nya ingin ditambahkan". \newline 3. Klik "Tambah Pod". \newline 4. Isi nama/nomor pod, pilih tipe pod, masukkan device code IoT yang unik. \newline 5. Klik "Simpan".          & - Pod baru berhasil ditambahkan. \newline - Status awal "Tersedia".                                                                     & Sesuai                    \\
    \hline
    Menambah pod baru dengan device code IoT yang sudah ada         & 1. Login. \newline 2. Ke "Manajemen Pod" - "Daftar Pod". \newline 3. Klik "Tambah Pod". \newline 4. Masukkan device code IoT yang sudah digunakan pod lain. \newline 5. Klik "Simpan".                                                              & - Pesan error "Device code sudah digunakan". \newline - Penambahan pod gagal.                                                           & Sesuai                    \\
    \hline
    Mengubah status pod menjadi "Perlu Dibersihkan" secara manual   & 1. Login. \newline 2. Ke "Manajemen Pod" - "Daftar Pod". \newline 3. Pilih pod yang "Tersedia", klik "Edit". \newline 4. Ubah status ke "Perlu Dibersihkan". \newline 5. Klik "Simpan".                                                             & - Status pod berhasil diubah.                                                                                                           & Sesuai                    \\
    \hline
    Membuat voucher diskon nominal untuk cabang                     & 1. Login. \newline 2. Ke "Manajemen Voucher Cabang". \newline 3. Klik "Tambah Voucher". \newline 4. Isi kode, nama, deskripsi, tipe "NOMINAL", jumlah diskon, min transaksi, periode berlaku, maks penggunaan. \newline 5. Klik "Simpan".           & - Voucher cabang berhasil dibuat.                                                                                                       & Sesuai                    \\
    \hline
    Membuat voucher diskon persentase dengan batas maks diskon      & 1. Login. \newline 2. Ke "Manajemen Voucher Cabang". \newline 3. Klik "Tambah Voucher". \newline 4. Isi kode, nama, tipe "PERSENTASE", nilai persentase, diskon maksimal, min transaksi, periode, kuota. \newline 5. Klik "Simpan".                 & - Voucher cabang berhasil dibuat.                                                                                                       & Sesuai                    \\
    \hline
    Menetapkan harga dinamis untuk akhir pekan                      & 1. Login. \newline 2. Ke "Harga Dinamis". \newline 3. Pilih tipe pod. \newline 4. Pilih rentang tanggal mencakup beberapa akhir pekan. \newline 5. Masukkan harga lebih tinggi. \newline  6. Klik "Simpan".                                         & - Harga dinamis untuk akhir pekan berhasil disimpan.                                                                                    & Sesuai                    \\
    \hline
    Menetapkan harga dinamis yang tumpang tindih dengan aturan lain & 1. Login. \newline 2. Ke "Harga Dinamis". \newline 3. Buat aturan harga dinamis baru yang tanggalnya tumpang tindih dengan aturan yang sudah ada untuk tipe pod yang sama.                                                                          & - Sistem memberikan peringatan bahwa aturan harga ada yang tumpang tindih dan aturan harga baru tidak berhasil disimmpan                & Sesuai                    \\
    \hline
    \multicolumn{4}{|l|}{\textbf{Fitur Manajemen Reservasi dan Tamu Cabang}}
    \\
    \hline
    Input manual transaksi OTA dengan semua field lengkap           & 1. Login. \newline 2. Ke "Transaksi OTA". \newline 3. Klik "Tambah". \newline 4. Isi Jenis OTA, ID Order OTA, nama tamu jika ada, etanggal check-in/out, tipe pod, jumlah tamu, dan harga di OTA \newline 5. Klik "Simpan".                         & - Transaksi OTA berhasil dicatat.                                                                                                       & Sesuai                    \\
    \hline
    Input manual transaksi OTA dengan field wajib kosong            & 1. Login. \newline 2. Ke "Transaksi OTA". \newline 3. Klik "Tambah". \newline 4. Kosongkan salah satu field wajib (misal: ID Order OTA). \newline 5. Coba klik "Simpan".                                                                            & - Tombol simpan tidak bisa diklik karena masih ada field wajib yang kosong                                                              & Sesuai                    \\
    \hline
    Memindahkan tamu ke pod lain karena ada masalah di pod asal     & 1. Login. \newline 2. Cari booking tamu aktif. \newline 3. Pilih "Pindah Pod", berikan alasan (misal: AC rusak). \newline 4. Pilih pod lain yang kosong, tipe sama. \newline 5. Konfirmasi.                                                         & - Tamu berhasil dipindahkan. \newline - Pod asal mungkin statusnya berubah jadi "Maintenance".                                          & Sesuai                    \\
    \hline
    Mencatat pembersihan pod dengan unggah foto bukti               & 1. Login. \newline 2. Ke "Log Kebersihan" atau daftar pod. \newline 3. Pilih pod "Perlu Dibersihkan". \newline 4. Klik "Sudah Dibersihkan". \newline 5. Unggah foto bukti kebersihan. \newline 6. (Jika ada) Tambahkan catatan. \newline 7. Simpan. & - Status pod "Tersedia". \newline - Foto bukti dan catatan tersimpan di log.                                                            & Sesuai                    \\
    \hline
    \multicolumn{4}{|l|}{\textbf{Fitur Manajemen Ulasan Cabang}}
    \\
    \hline
    Menyetujui ulasan pengguna yang konstruktif                     & 1. Login. \newline 2. Ke "Manajemen Ulasan". \newline 3. Pilih ulasan "Pending" yang isinya baik. \newline 4. Klik "Setujui/Tampilkan".                                                                                                             & - Ulasan tampil di halaman publik.                                                                                                      & Sesuai                    \\
    \hline
    Menolak ulasan pengguna yang mengandung spam/kata kasar         & 1. Login. \newline 2. Ke "Manajemen Ulasan". \newline 3. Pilih ulasan yang berisi spam. \newline 4. Klik "Tolak/Sembunyikan".                                                                                                                       & - Ulasan tidak tampil dan statusnya diperbarui.                                                                                         & Sesuai                    \\
    \hline
\end{longtable}


\begin{longtable}{|p{3.8cm}|p{4.5cm}|p{4.5cm}|p{2cm}|} % Menggunakan format 4 kolom
    \caption{Hasil Pengujian Black Box (SuperAdmin) } \label{tabel:hasil_bb_superadmin}                                                                                                                                                                                                                                                                                                                                                                                                                                                 \\
    \hline
    \textbf{Skenario Pengujian}                                    & \textbf{Langkah Pengujian}                                                                                                                                                                                                                                             & \textbf{Hasil yang diharapkan}                                                                                                                                & \textbf{Kesesuaian Hasil} \\
    \hline
    \endfirsthead

    % \multicolumn{4}{c}%
    % {{\bfseries Tabel \thetable\ lanjut dari halaman sebelumnya}}                                                                                                                                                                                                                                                                                                                                                                                                                                                                       \\
    % \hline
    % \textbf{Skenario Pengujian}                                    & \textbf{Langkah Pengujian}                                                                                                                                                                                                                                             & \textbf{Hasil yang diharapkan}                                                                                                                                & \textbf{Kesesuaian Hasil} \\
    % \hline
    % \endhead

    % \hline \multicolumn{4}{r}{\textit{Bersambung ke halaman berikutnya}}                                                                                                                                                                                                                                                                                                                                                                                                                                                                \\
    % \endfoot

    \hline
    \endlastfoot


    % --- Konten tabel Admin Pusat (SuperAdmin) ---
    \multicolumn{4}{|l|}{\textbf{Fitur Login \& Dashboard Global SuperAdmin}}
    \\
    \hline
    Login SuperAdmin berhasil                                      & 1. Buka halaman login admin. \newline 2. Masukkan email SuperAdmin. \newline 3. Masukkan kata sandi yang benar. \newline 4. Klik "Masuk".                                                                                                                              & - Login berhasil. \newline - Diarahkan ke dashboard SuperAdmin.                                                                                               & Sesuai                    \\
    \hline
    Melihat statistik global di dashboard SuperAdmin               & 1. Login sebagai SuperAdmin. \newline 2. Amati informasi di dashboard (total cabang, total pengguna, total pendapatan global, dll.).                                                                                                                                   & - Menampilkan data statistik agregat dari semua cabang yang akurat.                                                                                           & Sesuai                    \\
    \hline
    \multicolumn{4}{|l|}{\textbf{Fitur Manajemen Data Master Global}}
    \\
    \hline
    Menambah data cabang (hotel) baru                              & 1. Login sebagai SuperAdmin. \newline 2. Navigasi ke "Manajemen Cabang/Hotel". \newline 3. Klik "Tambah Cabang Baru". \newline 4. Isi semua field yang diperlukan (nama, alamat, kota, kontak, slug unik, dll.). \newline 5. Klik "Simpan".                            & - Cabang baru berhasil ditambahkan ke sistem. \newline - Cabang muncul di daftar cabang.                                                                      & Sesuai                    \\
    \hline
    Mengedit data cabang (hotel) yang sudah ada                    & 1. Login sebagai SuperAdmin. \newline 2. Navigasi ke "Manajemen Cabang/Hotel". \newline 3. Pilih cabang, klik "Edit". \newline 4. Ubah salah satu field (misal: email kontak). \newline 5. Klik "Simpan".                                                              & - Data cabang berhasil diperbarui.                                                                                                                            & Sesuai                    \\
    \hline
    Menambah data kota baru                                        & 1. Login sebagai SuperAdmin. \newline 2. Navigasi ke "Manajemen Kota". \newline 3. Klik "Tambah Kota". \newline 4. Masukkan nama kota. \newline 5. Klik "Simpan".                                                                                                      & - Kota baru berhasil ditambahkan. \newline - Kota dapat dipilih saat menambah/mengedit cabang.                                                                & Sesuai                    \\
    \hline
    Menambah akun Admin Cabang baru dan mengassign ke cabang       & 1. Login sebagai SuperAdmin. \newline 2. Navigasi ke "Manajemen Staff". \newline 3. Klik "Tambah Staff". \newline 4. Isi detail staff (nama, email, password), pilih peran "Admin Cabang". \newline 5. Assign ke satu atau beberapa cabang. \newline 6. Klik "Simpan". & - Akun Admin Cabang baru berhasil dibuat. \newline - Admin Cabang tersebut dapat login dan mengelola cabang yang di-assign.                                   & Sesuai                    \\
    \hline
    Menghapus akun staff (Admin Cabang)                            & 1. Login sebagai SuperAdmin. \newline 2. Navigasi ke "Manajemen Staff". \newline 3. Cari akun Admin Cabang. \newline 4. Klik opsi "Hapus".                                                                                                                             & - Akun Admin Cabang berhasil dihapus. \newline - Admin Cabang tersebut tidak dapat login.                                                                     & Sesuai                    \\
    \hline
    Menambah fasilitas global baru                                 & 1. Login sebagai SuperAdmin. \newline 2. Navigasi ke "Manajemen Fasilitas". \newline 3. Klik "Tambah Fasilitas". \newline 4. Masukkan nama fasilitas dan unggah ikon. \newline 5. Klik "Simpan".                                                                       & - Fasilitas baru berhasil ditambahkan. \newline - Fasilitas dapat di-assign ke hotel atau tipe pod.                                                           & Sesuai                    \\
    \hline
    Membuat voucher global untuk semua cabang                      & 1. Login. \newline 2. Ke "Manajemen Voucher". \newline 3. Klik "Tambah Voucher". \newline 4. Isi detail (kode, nama, tipe, nilai, periode, kuota). \newline 5. Pilih opsi "Berlaku di semua cabang". \newline 6. Klik "Simpan".                                        & - Voucher global berhasil dibuat. \newline - Dapat digunakan pengguna di semua cabang.                                                                        & Sesuai                    \\
    \hline
    \multicolumn{4}{|l|}{\textbf{Fitur Verifikasi \& Approval Sistem}}
    \\
    \hline
    Menyetujui permintaan verifikasi KYC pengguna yang valid       & 1. Login sebagai SuperAdmin. \newline 2. Navigasi ke "Verifikasi KYC". \newline 3. Pilih permintaan KYC yang berstatus "Pending" dengan data dan foto valid. \newline 4. Klik "Setujui".                                                                               & - Status KYC pengguna berubah menjadi "Approved". \newline - Pengguna mendapatkan pembaharuan. \newline - Pengguna dapat mengakses fitur yang memerlukan KYC. & Sesuai                    \\
    \hline
    Menolak permintaan verifikasi KYC pengguna (misal: foto buram) & 1. Login sebagai SuperAdmin. \newline 2. Navigasi ke "Verifikasi KYC". \newline 3. Pilih permintaan KYC yang fotonya buram. \newline 4. Klik "Tolak". \newline 5. Masukkan alasan penolakan (misal: "Foto tidak jelas").                                               & - Status KYC pengguna berubah menjadi "Rejected". \newline - Pengguna mendapatkan pembaharuan beserta alasan penolakan.                                       & Sesuai                    \\
    \hline
    Menyetujui permintaan pencairan saldo referral yang valid      & 1. Login sebagai SuperAdmin. \newline 2. Navigasi ke "Permintaan Payout Referral". \newline 3. Pilih permintaan yang valid (saldo pengguna cukup, data bank benar). \newline 4. Klik "Setujui". \newline 5. Proses transfer dana dan unggah buktinya.                  & - Status permintaan payout berubah menjadi "Approved". \newline - Pengguna mendapatkan pembaharuan.                                                           & Sesuai                    \\
    \hline
    Menolak permintaan pencairan saldo referral (saldo kurang)     & 1. Login sebagai SuperAdmin. \newline 2. Navigasi ke "Permintaan Payout Referral". \newline 3. Pilih permintaan dimana saldo aktual pengguna kurang dari yang diajukan. \newline 4. Klik "Tolak". \newline 5. Masukkan alasan.                                         & - Status permintaan payout berubah menjadi "Rejected". \newline - Pengguna mendapatkan pembaharuan.                                                           & Sesuai                    \\
    \hline
\end{longtable}


% \subsection{Pengujian Beban}

% Pengujian beban dilakukan

\subsection{\textit{User Acceptance Test}}

\textit{User Acceptance Test} (UAT) dilakukan untuk mengetahui apakah sistem yang telah dibuat dapat diterima oleh \textit{User}. UAT dilakukan dengan cara meminta \textit{User} untuk mencoba sistem dan memberikan feedback terhadap sistem yang telah dibuat, berikut pada tabel \ref{tabel:daftar_responden_uat} adalah daftar responden yang mengikuti UAT.


\begin{longtable}{|c|p{3.5cm}|p{3.5cm}|p{3.5cm}|c|}
    \caption{Daftar Responden \textit{User Acceptance Test} (UAT)}
    \label{tabel:daftar_responden_uat}
    \\
    \hline
    \textbf{No} & \textbf{Nama}               & \textbf{Peran (role)}                                  & \textbf{Profesi}     & \textbf{Umur} \\
    \hline
    \endfirsthead

    % \multicolumn{5}{c}%
    % {{\bfseries Tabel \thetable\ lanjut dari halaman sebelumnya}}
    % \\
    % \hline
    % \textbf{No} & \textbf{Nama}               & \textbf{Peran (role)}                           & \textbf{Profesi}     & \textbf{Umur} \\
    % \hline
    % \endhead

    % \hline \multicolumn{5}{r}{\textit{Bersambung ke halaman berikutnya}}                                                               \\
    % \endfoot

    \hline
    \endlastfoot

    1           & Reisika Jenny               & \textit{SuperAdmin} dan \newline \textit{Admin Cabang} & Business Development & 22            \\ \hline
    2           & Asraf Rafli Daifuloh        & \textit{User}                                          & Fullstack Programmer & 24            \\ \hline
    3           & Damianus Cahyo Widi Nugroho & \textit{User}                                          & AI Engineer          & 22            \\ \hline
    4           & Evi Nur Hidayah             & \textit{User}                                          & Mahasiwa             & 21            \\ \hline
    5           & Moh. Eka Syafrino Nazhifan  & \textit{User}                                          & Fullstack Programmer & 19            \\ \hline
    6           & Salsa Ainnur Muharramah     & \textit{User}                                          & Mahasiswa Profesi    & 22            \\ \hline
    7           & Daffa Haj Tsaqif            & \textit{User}                                          & IoT Engineer         & 25            \\ \hline
    8           & Markus Rabin Ronaldo        & \textit{User}                                          & Mahasiswa            & 23            \\ \hline
    9           & Bela Dwi Primasari          & \textit{User}                                          & Sistem Analis        & 23            \\ \hline
    10          & Nazil Makarim Ramadhany     & \textit{User}                                          & Graphic Designer     & 21            \\ \hline
    11          & Bagus erwanto               & \textit{User}                                          & IT Support           & 23            \\ \hline
\end{longtable}

Untuk mendapatkan hasil yang lebih terukur dan seimbang setelah para responden mencoba sistem, penulis memberikan kuesioner yang berisi 10 pertanyaan yang harus dijawab oleh para responden. Kuesioner ini bertujuan untuk mengetahui sejauh mana sistem yang telah dibuat dapat memenuhi kebutuhan \textit{User}.

\begin{longtable}{|p{0.8cm}|p{7.5cm}|>{\centering\arraybackslash}p{0.8cm}|>{\centering\arraybackslash}p{0.8cm}|>{\centering\arraybackslash}p{0.8cm}|>{\centering\arraybackslash}p{0.8cm}|>{\centering\arraybackslash}p{0.8cm}|}
    \caption{Hasil Kuesioner \textit{User Acceptance Test} - Perspektif \textit{User}} \label{tabel:uat_user_hasil}                                                                                                                                   \\
    \hline
    \multirow{2}{*}{\textbf{No.}}              & \multirow{2}{*}{\textbf{Pertanyaan}}                                                                   & \multicolumn{5}{c|}{\textbf{Jawaban}}                                                       \\ \cline{3-7}
                                               &                                                                                                        & \textbf{STS}                          & \textbf{TS} & \textbf{N} & \textbf{S} & \textbf{SS} \\
    \hline
    \endfirsthead

    \multicolumn{7}{c}%
    {{\bfseries Tabel \thetable\ lanjut dari halaman sebelumnya}}                                                                                                                                                                                     \\
    \hline
    \multirow{2}{*}{\textbf{No.}}              & \multirow{2}{*}{\textbf{Pertanyaan}}                                                                   & \multicolumn{5}{c|}{\textbf{Jawaban}}                                                       \\ \cline{3-7}
                                               &                                                                                                        & \textbf{STS}                          & \textbf{TS} & \textbf{N} & \textbf{S} & \textbf{SS} \\
    \hline
    \endhead

    \hline \multicolumn{7}{r}{\textit{Bersambung ke halaman berikutnya}}                                                                                                                                                                              \\
    \endfoot

    \hline % \endlastfoot akan menjadi baris terakhir tabel setelah "Total Skor"
    \endlastfoot

    % --- Pertanyaan UAT untuk Pengguna (10 Pertanyaan) ---
    1.                                         & Apakah proses pendaftaran akun dan login ke aplikasi Tamu.id mudah dipahami dan dilakukan?             &                                       &             &            & 1          & 9           \\
    \hline
    2.                                         & Apakah tampilan antarmuka aplikasi Tamu.id menarik dan mudah digunakan secara umum?                    &                                       &             & 1          & 3          & 6           \\
    \hline
    3.                                         & Apakah informasi mengenai cabang dan tipe pod (harga, fasilitas, ketersediaan) disajikan dengan jelas? &                                       &             &            & 3          & 7           \\
    \hline
    4.                                         & Apakah alur proses pencarian dan pemesanan pod mudah diikuti dan efisien?                              &                                       &             &            & 3          & 7           \\
    \hline
    5.                                         & Apakah proses pembayaran saat melakukan pemesanan mudah dan terasa aman?                               &                                       &             & 1          & 1          & 8           \\
    \hline
    6.                                         & Apakah proses check-in dan check-out melalui aplikasi berjalan dengan lancar?                          &                                       &             &            &            & 10          \\
    \hline
    7.                                         & Apakah penggunaan QR Code untuk akses pod (jika ada fitur ini) mudah dan praktis?                      &                                       &             &            & 1          & 9           \\
    \hline
    8.                                         & Apakah saya mudah menemukan riwayat pemesanan dan detail transaksi saya?                               &                                       &             & 1          & 3          & 6           \\
    \hline
    9.                                         & Apakah fitur untuk memberikan ulasan dan rating setelah menginap mudah digunakan?                      &                                       &             &            & 5          & 5           \\
    \hline
    10.                                        & Secara keseluruhan, apakah Anda puas dengan fungsionalitas dan kemudahan penggunaan aplikasi Tamu.id?  &                                       &             &            & 1          & 9           \\
    \hline
    % --- Baris Total Bobot, Skor, Total Skor (Format Baru) ---
    \multicolumn{2}{|l|}{\textbf{Total Bobot}} &                                                                                                        &                                       & 3           & 21         & 76                       \\
    \hline
    \multicolumn{2}{|l|}{\textbf{Skor}}        &                                                                                                        &                                       & 9           & 84         & 380                      \\
    \hline
    \multicolumn{2}{|l|}{\textbf{Total Skor}}  & \multicolumn{5}{c|}{473}                                                                                                                                                                             \\ % Mengganti "anjay" dengan "..."
\end{longtable}

Dari tabel \ref{tabel:uat_user_hasil} di atas, kita dapat menghitung skor \textit{User Acceptance Test} (UAT) untuk perspektif pengguna sebagai berikut:

\begin{align*}
    \text{Skor UAT User}(\%) & = \left( \frac{A}{B \times N} \right) \times 100\%     \\
                             & = \left( \frac{473}{50 \times 10} \right) \times 100\% \\
                             & = \left( \frac{473}{500} \right) \times 100\%          \\
                             & = 0.946 \times 100\%                                   \\
                             & = 94.6\%
\end{align*}

Kesimpulan dari hasil UAT perspektif pengguna menunjukkan bahwa aplikasi Tamu.id mendapatkan skor 94.6\%, yang berarti aplikasi ini sangat diterima oleh pengguna dan masuk dalam kategori Sangat Layak berdasaarkan tabel \ref{tabel:kategori-kelayakan}.

\begin{longtable}{|p{0.8cm}|p{7.5cm}|>{\centering\arraybackslash}p{0.8cm}|>{\centering\arraybackslash}p{0.8cm}|>{\centering\arraybackslash}p{0.8cm}|>{\centering\arraybackslash}p{0.8cm}|>{\centering\arraybackslash}p{0.8cm}|}
    \caption{Hasil Kuesioner \textit{User Acceptance Test} - Perspektif \textit{Admin} Cabang} \label{tabel:uat_admin_cabang_hasil}                                                                                                                                                    \\
    \hline
    \multirow{2}{*}{\textbf{No.}}              & \multirow{2}{*}{\textbf{Pertanyaan}}                                                                                                    & \multicolumn{5}{c|}{\textbf{Jawaban}}                                                       \\ \cline{3-7}
                                               &                                                                                                                                         & \textbf{STS}                          & \textbf{TS} & \textbf{N} & \textbf{S} & \textbf{SS} \\
    \hline
    \endfirsthead

    \multicolumn{7}{c}%
    {{\bfseries Tabel \thetable\ lanjut dari halaman sebelumnya}}                                                                                                                                                                                                                      \\
    \hline
    \multirow{2}{*}{\textbf{No.}}              & \multirow{2}{*}{\textbf{Pertanyaan}}                                                                                                    & \multicolumn{5}{c|}{\textbf{Jawaban}}                                                       \\ \cline{3-7}
                                               &                                                                                                                                         & \textbf{STS}                          & \textbf{TS} & \textbf{N} & \textbf{S} & \textbf{SS} \\
    \hline
    \endhead

    \hline \multicolumn{7}{r}{\textit{Bersambung ke halaman berikutnya}}                                                                                                                                                                                                               \\
    \endfoot

    \hline % \endlastfoot
    \endlastfoot

    % --- Pertanyaan UAT untuk Admin Cabang (10 Pertanyaan) ---
    1.                                         & Apakah proses login ke dashboard Admin Cabang mudah dan aman?                                                                           &                                       &             &            &            & 1           \\
    \hline
    2.                                         & Apakah tampilan dashboard Admin Cabang informatif untuk memantau operasional harian cabang?                                             &                                       &             &            &            & 1           \\
    \hline
    3.                                         & Apakah pengelolaan informasi umum cabang (seperti detail kontak, foto, atau deskripsi) mudah dilakukan?                                 &                                       &             &            &            & 1           \\
    \hline
    4.                                         & Apakah fitur untuk menambah dan mengelola tipe pod serta unit pod individual di cabang saya mudah digunakan?                            &                                       &             &            &            & 1           \\
    \hline
    5.                                         & Apakah proses untuk mengatur harga pod (termasuk harga dinamis jika digunakan) cukup jelas dan mudah diterapkan?                        &                                       &             &            &            & 1           \\
    \hline
    6.                                         & Apakah pengelolaan data reservasi tamu (melihat daftar booking, detail, dan status) efisien?                                            &                                       &             &            &            & 1           \\
    \hline
    7.                                         & Apakah fitur untuk menginput atau mengelola data reservasi yang berasal dari Online Travel Agent (OTA) mudah dioperasikan?              &                                       &             &            &            & 1           \\
    \hline
    8.                                         & Apakah proses untuk mengelola status kebersihan dan ketersediaan pod (misalnya, dari "Perlu Dibersihkan" menjadi "Tersedia") sederhana? &                                       &             &            &            & 1           \\
    \hline
    9.                                         & Apakah fitur untuk melihat dan menanggapi (memoderasi) ulasan dari pengguna untuk cabang saya mudah digunakan?                          &                                       &             &            &            & 1           \\
    \hline
    10.                                        & Secara keseluruhan, apakah dashboard Admin Cabang ini efektif membantu saya dalam mengelola operasional cabang sehari-hari?             &                                       &             &            &            & 1           \\
    \hline
    % --- Baris Total Bobot, Skor, Total Skor (Format Baru) ---
    \multicolumn{2}{|l|}{\textbf{Total Bobot}} &                                                                                                                                         &                                       &             &            & 10                       \\
    \hline
    \multicolumn{2}{|l|}{\textbf{Skor}}        &                                                                                                                                         &                                       &             &            & 50                       \\
    \hline
    \multicolumn{2}{|l|}{\textbf{Total Skor}}  & \multicolumn{5}{c|}{50}                                                                                                                                                                                                               \\ % Mengganti "anjay" dengan "..."
\end{longtable}

Dari tabel \ref{tabel:uat_admin_cabang_hasil} di atas, kita dapat menghitung skor \textit{User Acceptance Test} (UAT) untuk perspektif Admin Cabang sebagai berikut:

\begin{align*}
    \text{Skor UAT Admin}(\%) & = \left( \frac{A}{B \times N} \right) \times 100\%   \\
                              & = \left( \frac{50}{50 \times 1} \right) \times 100\% \\
                              & = \left( \frac{50}{50} \right) \times 100\%          \\
                              & = 1 \times 100\%                                     \\
                              & = 100\%
\end{align*}

Kesimpulan dari hasil UAT perspektif Admin Cabang menunjukkan bahwa aplikasi Tamu.id mendapatkan skor 100\%, yang berarti aplikasi ini sangat diterima oleh Admin Cabang dan masuk dalam kategori Sangat Layak berdasar tabel \ref{tabel:kategori-kelayakan}.


\begin{longtable}{|p{0.8cm}|p{7.5cm}|>{\centering\arraybackslash}p{0.8cm}|>{\centering\arraybackslash}p{0.8cm}|>{\centering\arraybackslash}p{0.8cm}|>{\centering\arraybackslash}p{0.8cm}|>{\centering\arraybackslash}p{0.8cm}|}
    \caption{Hasil Kuesioner \textit{User Acceptance Test} - Perspektif \textit{SuperAdmin}} \label{tabel:uat_superadmin}                                                                                                                                                               \\
    \hline
    \multirow{2}{*}{\textbf{No.}}              & \multirow{2}{*}{\textbf{Pertanyaan}}                                                                                                     & \multicolumn{5}{c|}{\textbf{Jawaban}}                                                       \\ \cline{3-7}
                                               &                                                                                                                                          & \textbf{STS}                          & \textbf{TS} & \textbf{N} & \textbf{S} & \textbf{SS} \\
    \hline
    \endfirsthead

    \multicolumn{7}{c}%
    {{\bfseries Tabel \thetable\ lanjut dari halaman sebelumnya}}                                                                                                                                                                                                                       \\
    \hline
    \multirow{2}{*}{\textbf{No.}}              & \multirow{2}{*}{\textbf{Pertanyaan}}                                                                                                     & \multicolumn{5}{c|}{\textbf{Jawaban}}                                                       \\ \cline{3-7}
                                               &                                                                                                                                          & \textbf{STS}                          & \textbf{TS} & \textbf{N} & \textbf{S} & \textbf{SS} \\
    \hline
    \endhead

    \hline \multicolumn{7}{r}{\textit{Bersambung ke halaman berikutnya}}                                                                                                                                                                                                                \\
    \endfoot

    \hline % \endlastfoot
    \endlastfoot

    % --- Pertanyaan UAT untuk SuperAdmin (10 Pertanyaan) ---
    1.                                         & Apakah proses login ke dashboard SuperAdmin mudah dan aman?                                                                              &                                       &             &            &            & 1           \\
    \hline
    2.                                         & Apakah tampilan dashboard global SuperAdmin menyajikan informasi ringkasan sistem secara efektif dan mudah dipahami?                     &                                       &             &            &            & 1           \\
    \hline
    3.                                         & Apakah pengelolaan data master sistem (seperti daftar Cabang/Hotel, Kota, OTA global) mudah dilakukan?                                   &                                       &             &            &            & 1           \\
    \hline
    4.                                         & Apakah fitur untuk menambah, mengedit, dan mengelola akun staf (Admin Cabang, SuperAdmin lain) serta hak aksesnya intuitif?              &                                       &             &            &            & 1           \\
    \hline
    5.                                         & Apakah proses verifikasi identitas pengguna (KYC) dari sisi SuperAdmin berjalan dengan alur yang jelas dan efisien?                      &                                       &             &            & 1          &             \\
    \hline
    6.                                         & Apakah proses approval atau penolakan permintaan pencairan saldo referral pengguna mudah dikelola dan dilacak?                           &                                       &             &            & 1          &             \\
    \hline
    7.                                         & Apakah fitur untuk membuat dan mengelola voucher yang berlaku secara global atau untuk banyak cabang mudah digunakan?                    &                                       &             &            &            & 1           \\
    \hline
    8.                                         & Apakah sistem SuperAdmin menyediakan alat yang memadai untuk memantau aktivitas penting dan kesehatan sistem secara keseluruhan?         &                                       &             &            &            & 1           \\
    \hline
    9.                                         & Apakah laporan global atau data analitik yang dihasilkan sistem (jika ada) membantu dalam analisis dan pengambilan keputusan strategis?  &                                       &             &            &            & 1           \\
    \hline
    10.                                        & Secara keseluruhan, apakah antarmuka SuperAdmin efektif dalam membantu saya menjalankan tugas administrasi sistem dan pengawasan global? &                                       &             &            &            & 1           \\
    \hline
    % --- Baris Total Bobot, Skor, Total Skor (Format Baru) ---
    \multicolumn{2}{|l|}{\textbf{Total Bobot}} &                                                                                                                                          &                                       &             & 2          & 8                        \\
    \hline
    \multicolumn{2}{|l|}{\textbf{Skor}}        &                                                                                                                                          &                                       &             & 8          & 40                       \\
    \hline
    \multicolumn{2}{|l|}{\textbf{Total Skor}}  & \multicolumn{5}{c|}{48}                                                                                                                                                                                                                \\ % Mengganti "anjay" dengan "..."
\end{longtable}

Dari tabel \ref{tabel:uat_superadmin} di atas, kita dapat menghitung skor \textit{User Acceptance Test} (UAT) untuk perspektif SuperAdmin sebagai berikut:

\begin{align*}
    \text{Skor UAT SuperAdmin}(\%) & = \left( \frac{A}{B \times N} \right) \times 100\%   \\
                                   & = \left( \frac{48}{50 \times 1} \right) \times 100\% \\
                                   & = \left( \frac{48}{50} \right) \times 100\%          \\
                                   & = 0.96 \times 100\%                                  \\
                                   & = 96\%
\end{align*}

Kesimpulan dari hasil UAT perspektif SuperAdmin menunjukkan bahwa aplikasi Tamu.id mendapatkan skor 96\%, yang berarti aplikasi ini sangat diterima oleh SuperAdmin dan masuk dalam kategori Sangat Layak berdasar tabel \ref{tabel:kategori-kelayakan}.